\documentclass[main]{subfiles}
%\newcommand{\materia}{ Máquinas Eléctricas II }
\newcommand{\curso}{6°1°}
\newcommand{\nombre}{Gustavo David Ferreyra}
\newcommand{\recursos}{\rotatebox[origin=r]{90}{Pizarrón, proyectores, videos, simuladores, libros de texto, materiales manipulativos y plataformas digitales.}}
\newcommand{\tiempos}{\rotatebox[origin=r]{90}{8 Semanas}}
\newcommand{\evaluacion}{Para evaluar el aprendizaje en relación a los saberes prioritarios del contenido, utilizan diferentes estrategias de evaluación, tales como pruebas escritas, ejercicios prácticos, proyectos individuales o en grupo, y observación del desempeño de los estudiantes. Es importante que estas estrategias estén alineadas con las estrategias de enseñanza utilizadas y que se proporcione retroalimentación constructiva a los estudiantes. Se debe evaluar tanto el conocimiento como la comprensión y la aplicación de los conceptos, así como las habilidades prácticas de los estudiantes.}
\begin{document}
\chapter*{Planificación de \materia }
\info{\materia}{

ESPACIO CURRICULAR:\materia

CURSO:\curso

PROFESOR/A: \nombre

TITULO DEL EGRESADO: Técnico en Electricidad
}

\subsection{ACUERDOS TRANSVERSALES INSTITUCIONALES:}
\begin{itemize}
\item  Abordar de forma transversal la comprensión lectora, fluidez, hábitos de lectura y normativa ortográfica.
\item Promover la comunicación informal entre docentes y coordinadores de área mediante el uso del correo electrónico, what´s app informativo.
\item Implementar diversos instrumentos de evaluación, con énfasis en la utilización de las TIC.
\item Incentivar la modalidad técnica de la institución a partir de la participación en programas y proyectos específicos del área.
\item Promover el trabajo interdisciplinario mediante la puesta en marcha de proyectos comunes, con la metodología de aprendizaje basado en proyectos .Estimular la implementación de
evaluaciones integrando los contenidos de diferentes espacios curriculares.
\item  Poner énfasis el logro de un aprendizaje significativo mediante la indagación respecto a conocimientos previos.
\item Evaluar a los alumnos de forma integral, formativa y continua, registrando periódicamente en el GEM y en las planillas correspondientes las calificaciones obtenidas por el alumno en su
proceso de aprendizaje
\end{itemize}
\subsection{ACUERDOS DEL ÁREA TÉCNICA 2023}

El presente acuerdo del Área Técnica estará incluido en las planificaciones correspondientes y representa el compromiso de cada Docente para el
presente ciclo lectivo.
\begin{itemize}
    \item Respetar y hacer cumplir las normas de convivencia de la Escuela.
\item Cuidar y controlar que los alumnos no dañen materiales, herramientas, bibliografía,
mobiliario y todo elemento de la Escuela.
\item Ayudar en el control del uso de la vestimenta permitida.
\item Controlar el uso del celular, tanto alumnos como Docentes.
\item Como contenido transversal de la escuela, en cada materia se realizará lectura de
textos y comprensión de los mismos. La exposición oral de los alumnos es fundamental
que se realice en cada año y es un compromiso acordado por los docentes.
\item Realizar la corrección de ortografía en todos los espacios curriculares.
\item Debido a la organización para el uso del laboratorio, los Docentes deben avisar con
una semana de anticipación que lo necesitan para prever el traslado de los alumnos
que diariamente lo ocupan hacia otro lugar de la Escuela. Existirá una planilla para que
soliciten su uso. Se cuenta con un armario móvil que contiene instrumentos para
realizar mediciones eléctricas.
\item En las planificaciones indicar los proyectos a realizar durante el año.
\item En la planificación debe estar principalmente el programa anual de la materia.
\item Las carpetas que se realicen en hoja normalizada, deben utilizar la nueva hoja que se
puede adquirir en biblioteca. El uso de la misma es debido a que existen distintos
formatos y principalmente para abaratar costos al alumno.
\item Para los alumnos de sexto año se adjuntará las incumbencias que corresponde al título
otorgado por la escuela.
\item Los proyectos finales serán comenzados en quinto año, en la selección de los mismos,
los docentes encargados del proyecto definirán la viabilidad de estos para realizarse
prácticamente en sexto año.
\item Se respetarán las fechas para colocar las notas correspondientes en el GEM.
\item Participar activamente en la ExpoMerín y fomentar la importancia de la misma a
nuestros alumnos.
\item Será una condición necesaria la presentación de carpetas cuando los alumnos
concurran a rendir una materia, la presentación de la misma tiene un puntaje a tener
en cuenta en dicho examen.
\item Si en alguna mesa de examen el Docente no pudiera concurrir, dejar la evaluación
correspondiente en dirección
\end{itemize}
\newpage
\small
\begin{longtable}{|p{0.5cm}|p{1.5cm}|p{3.5cm}|p{3.5cm}|p{4cm}|p{.5cm}|p{0.5cm}|} \hline
   \rotatebox[origin=r]{90}{EJE} & SABERES & APRENDIZAJES & ESTRATEGIAS & EVALUACIÓN & \rotatebox[origin=r]{90}{RECURSOS} & \rotatebox[origin=r]{90}{TIEMPOS} \\ \hline \hline \endhead
    \rotatebox[origin=r]{90}{FUNDAMENTOS DE MÁQUINAS MONOFÁSICAS} & Reconocer los principios
básicos que fundamentan el
funcionamiento de una
máquina eléctrica con
energía alterna.& \begin{itemize}
\item Explicitación de las leyes electromagnéticas presentes en
una máquina eléctrica.
\item Interpretación de los campos electromagnéticos
originados por corriente alterna.
\item Distinción de la frecuencia y la velocidad de sincronismo
en una máquina eléctrica.
\item  Caracterización de la electrodinámica en máquinas
industriales. \end{itemize} & \begin{itemize}
    \item Utilizar modelos físicos o virtuales para mostrar los componentes de un motor asíncrono y cómo se relacionan entre sí.
\item Proporcionar imágenes y diagramas de motores asíncronos para que los estudiantes puedan identificar las diferentes partes y su función.
\item Utilizar vídeos que muestren la construcción y el ensamblaje de un motor asíncrono para ayudar a los estudiantes a comprender mejor cómo funciona.
\end{itemize}&  \evaluacion & \recursos & \tiempos \\ \hline

& Interpretar las teorías
electrodinámicas eléctricas
en una máquina rotativa. & \begin{itemize}
    \item Interpretación de los ángulos geométricos y eléctricos.
\item Comprensión de la fuerza electromotriz inducida en
máquinas rotativas.
\item Reconocimiento de los factores de arrollamiento y de
paso.
\item Análisis enmarcado en la tecnológica eléctrica de los
usos y características de los motores paso a paso. 
\end{itemize}& Para enseñar la interpretación de ángulos geométricos y eléctricos, se pueden utilizar herramientas visuales y experimentales para que los estudiantes comprendan su aplicación en la resolución de problemas. En cuanto a la comprensión de la fuerza electromotriz inducida en máquinas rotativas, es importante enseñar los conceptos fundamentales de la inducción electromagnética y utilizar ejemplos concretos para ayudar a los estudiantes a comprender su aplicación práctica. El reconocimiento de factores de arrollamiento y de paso se puede enseñar mediante ejemplos que muestren cómo afectan a la resistencia y la conductividad en los circuitos eléctricos. Finalmente, el análisis de motores paso a paso debe centrarse en su uso y características en la tecnología eléctrica, y se puede complementar con demostraciones y ejemplos prácticos.&  \evaluacion & \recursos & \tiempos \\ \hline


% eje 2
\rotatebox[origin=r]{90}{MOTORES ASINCRÓNICOS TRIFÁSICOS EN DIVERSOS ENSAYOS ELÉCTRICOS} & Conocer la naturaleza
eléctrica de las máquinas
asincrónicas. &\begin{itemize} \item Interpretación de los aspectos constructivos de un motor
asincrónico.
\item  Reconocimiento del principio de funcionamiento de una
máquina asíncrona y su resbalamiento.
\item Identificación del motor eléctrico como transformador
energético en su empleo técnico.
\item  Análisis del comportamiento eléctrico con rotor detenido y
en marcha. \end{itemize} 
& Se pueden utilizar imágenes, videos y simuladores para visualizar el funcionamiento del motor asíncrono y sus partes. Además, fomentar la discusión en el aula sobre los diferentes aspectos del motor y su aplicación en la vida cotidiana. Actividades prácticas como mediciones, ensamblaje y cálculo del resbalamiento del motor también son útiles para aplicar los conocimientos adquiridos. &\evaluacion & \recursos & \tiempos \\ \hline

& Analizar críticamente los
ensayos eléctricos
característicos de un motor
asincrónico que
fundamentan los distintos
fenómenos energéticos.& \begin{itemize}
    \item Elaboración de los diagramas vectoriales.
\item Interpretación gráfica de los circuitos equivalentes
técnicamente eléctricos.
\item Análisis de las curvas de arranque, cuplas y balance de
potencias.
\item Determinación de diagramas, curvas y comportamientos
característicos de las máquinas asincrónicas en ensayos
eléctricos en el Laboratorio.
\item Distinción de formas de arranque de un motor trifásico.
\item Reconocimiento de la máquina asíncrona como
autotransformador.
\item Contrastación de los distintos tipos de frenado eléctrico
en sus distintas concepciones tecnológicas. 

\end{itemize}& Para enseñar la elaboración de diagramas vectoriales, la interpretación gráfica de circuitos equivalentes eléctricos, análisis de curvas de arranque, cuadros y balance de potencias, determinación de comportamientos característicos de las máquinas asíncronas en ensayos eléctricos en el laboratorio, distinción de formas de arranque de un motor trifásico, reconocimiento de la máquina asíncrona como autotransformador y contrastación de distintos tipos de frenado eléctrico, se pueden utilizar estrategias como la utilización de simulaciones, ejemplos prácticos y la realización de actividades en el laboratorio. Esto permite a los estudiantes experimentar y comprender de manera más efectiva estos conceptos complejos. & \evaluacion & \recursos & \tiempos \\ \hline


\rotatebox[origin=r]{90}{ENSAYOS EN MÁQUINAS SINCRÓNICAS} & Comprender el
funcionamiento de las
partes constitutivas de una
máquina sincrónica. &\begin{itemize}
\item Identificación de los principios de funcionamiento del
generador y del motor eléctrico.
\item Descripción de las máquinas sincrónicas en vacío y en
carga.
\item Interpretación de las partes constitutivas de las máquinas
sincrónicas.
\item Reformulación de la frecuencia de la fuerza electromotriz
inducida.
\item Caracterización de la reacción de inducido.
\item Diferenciación en la regulación de una máquina
sincrónica. 
\end{itemize} & Para enseñar la identificación de los principios de funcionamiento del generador y del motor eléctrico, descripción de las máquinas sincrónicas en vacío y en carga, interpretación de las partes constitutivas de las máquinas sincrónicas, reformulación de la frecuencia de la fuerza electromotriz inducida, caracterización de la reacción de inducido y diferenciación en la regulación de una máquina sincrónica, se pueden utilizar estrategias como la realización de ejemplos prácticos, la utilización de videos y simulaciones, la discusión en el aula y la resolución de problemas. Además, actividades en el laboratorio pueden ser muy útiles para experimentar y comprender estos conceptos complejos. Esto permite a los estudiantes aplicar los conocimientos teóricos y adquirir habilidades prácticas importantes para su futuro desempeño profesional.& \evaluacion & \recursos& \tiempos \\ \hline

 & Valorar la utilidad industrial
de una máquina sincrónica
considerando su
importancia tecnológica
eléctrica. & \begin{itemize}
    \item Selección de la metodología de sincronización de
generadores sincrónicos.
\item Caracterización de una puesta en paralelo de
generadores de aplicación industrial eléctrica.
\item Indagación del uso de una máquina sincrónica en la
corrección del factor de potencia.
\item Análisis de las características técnicas de una máquina
sincrónica, mediante el empleo de las metodologías y los
ensayos en el Laboratorio Eléctrico.
\item Reconocimiento tecnológico de una máquina sincrónica
en su rol de generador o motor eléctrico.
\item Análisis crítico de la fuerza contraelectromotriz y los
principios de funcionamiento de un motor sincrónico. 
\end{itemize}& Para enseñar la selección de la metodología de sincronización de generadores sincrónicos, se pueden utilizar ejemplos prácticos y experiencias simuladas que permitan a los estudiantes entender las ventajas y desventajas de cada metodología. En cuanto a la puesta en paralelo de generadores, es importante enfatizar la seguridad y la confiabilidad en la operación de sistemas eléctricos de gran escala. Para comprender el uso de una máquina sincrónica en la corrección del factor de potencia, se pueden utilizar demostraciones experimentales y simulaciones que permitan visualizar la aplicación práctica. Enseñar el análisis crítico de la fuerza contraelectromotriz y los principios de funcionamiento de un motor sincrónico, se pueden emplear recursos audiovisuales y herramientas de simulación que permitan a los estudiantes comprender de manera más visual y práctica los conceptos teóricos.& \evaluacion & \recursos& \tiempos \\ \hline




\end{longtable}

\end{document}
